\section{The Model}
\label{sec:model}


The model follows households as they receive income, spend, borrow and repay debt.
Each household begins with financial characteristics drawn from survey microdata.
Traditional credit is subject to a residual-income test based on Regulation 23A,
while \ac{BNPL} platforms apply a lighter screen.

This section explains who acts in the model, what each institution can observe,
and how household finances change over time. Parameter choices and calibration
are discussed in Section~\ref{sec:calibration}. The specification follows the
\ac{ODD} protocol \cite{grimm2006odd, grimm2020odd}, with the operational rules
summarised in Table~\ref{tab:submodels}. Appendix~\ref{app:specification} records
the design concepts and exclusions, Appendix~\ref{app:rationale} explains the
alternatives considered, and Appendices~\ref{app:pseudo-init}
and~\ref{app:pseudo-main} provide the pseudocode.

\label{sec:purpose}
The experiment begins with a traditional-credit baseline calibrated to 2017
conditions. We then enable \ac{BNPL} while holding the other parameters fixed.
All households start with zero \ac{BNPL} balances, so any subsequent exposure is
generated during the simulation. Comparing the runs isolates the effect of the
additional credit channel within the model.

\subsection{Agents and Institutions}
\label{sec:entities}

The model contains five types of entity:
\begin{itemize}
    \item \textbf{Households} respond to their current finances through the borrowing
    and repayment rules below. Each is based on a \ac{NIDS} Wave 5 household
    (Table~\ref{tab:agent-mapping}) and tracks savings, traditional debt,
    \ac{BNPL} obligations, arrears, distress and default. Its repayment type,
    minimum-payer or scheduled-payer, is fixed at initialisation.
    \item \textbf{The traditional lender} holds all traditional debt in one
    consolidated institution. It applies a fixed affordability rule and does not
    adapt its lending policy in response to losses.
    \item \textbf{The credit bureau} records traditional debt and defaults.
    Its contents determine which recorded obligations enter the traditional
    lender's assessment.
    \item \textbf{\ac{BNPL} platforms} provide the additional credit channel.
    The benchmark uses four platforms. Each maintains its own receivables,
    applies an automated screen, and limits both individual orders and each
    household's outstanding exposure. Platforms cannot observe competitors'
    balances.
    \item \textbf{Reference groups} define whose adoption a household observes.
    Each household belongs to the group formed by its income quintile and province.
\end{itemize}

Table~\ref{tab:agent-mapping} links the initial household variables to their
sources. These include observed survey fields, imputed characteristics and
constructed financial quantities. Section~\ref{sec:calibration} explains their
construction.

\begin{table}[H]
\centering
\caption{Data-to-agent mapping}
\label{tab:agent-mapping}
\tabnote{The state variables initialised for each agent, their meaning, and their source. Monetary fields are 2017 Rands. Fields marked \emph{tick} are monthly survey flows scaled by $12/26$ to the 14-day tick; stocks (savings, debt) are not scaled. Construction detail for every derived field is in Appendix~\ref{app:variables}; the imputed flags are described in Section~\ref{sec:match} and Appendix~\ref{app:matching}. Two dimensions are not observed in the microdata: every agent starts with a zero \ac{BNPL} balance, which is the injection design, and the repayment archetype of Section~\ref{sec:submodels} is assigned at random at initialisation.}
\begin{tabular}{L{4.8cm}L{4.4cm}L{4.2cm}}
\toprule
\textbf{Agent field} & \textbf{Description} & \textbf{Source} \\
\midrule
\texttt{income\_monthly} & Total monthly liquidity inflows & NIDS \\
\texttt{income\_wage\_tick} & Wage income exposed to shocks & NIDS \\
\texttt{n\_earners} & Employed members on the roster & NIDS individual file \\
\texttt{income\_source} & Dominant component: wage, grant, other & NIDS, derived \\
\texttt{committed\_tick} & Incompressible food and rent costs & NIDS, derived \\
\texttt{discretionary\_tick} & Non-food expenditure & NIDS, derived \\
\texttt{liquid\_savings} & Proxy for immediate cash buffer & NIDS, winsorised \\
\texttt{d\_trad} & Outstanding traditional debt balance & NIDS \\
\texttt{apr\_annual}, \texttt{term\_months} & Rate and horizon pricing that balance & Product mix (Section~\ref{sec:servicing}) \\
\texttt{scheduled\_service\_tick} & Debt servicing obligation per tick & Constructed (Section~\ref{sec:servicing}) \\
\texttt{nca\_capacity\_monthly} & Regulation 23A affordable residual & Statutory (Section~\ref{sec:servicing}) \\
\texttt{banked} & Eligibility gate for \ac{BNPL} & FinScope, imputed \\
\texttt{income\_quintile} & Income quintile; match key & NIDS, weighted \\
\texttt{province} & Geographic tag; donor match key & NIDS \\
\texttt{reference\_group} & Group for peer adoption & Derived (quintile $\times$ province) \\
\texttt{household\_size} & Total household members & NIDS \\
\bottomrule
\end{tabular}
\end{table}

\subsection{Balance Sheets and Information}
\label{sec:asymmetry}

Household finances are tracked through income, expenditure, savings and debt.
Food and rent form committed expenditure, which the household cannot reduce.
Discretionary spending can fall to zero to make room for repayments, and liquid
savings provide a buffer when income is insufficient. Traditional debt and
\ac{BNPL} balances are recorded separately because the lenders have different
information about them.

Figure~\ref{fig:visibility} shows this information structure. In the benchmark,
platforms cannot see competitors' exposures and the bureau does not record
\ac{BNPL}, reflecting the reporting position described in
Section~\ref{sec:background}.

\begin{figure}[H]
\centering
\begin{tikzpicture}[
    font=\small,
    box/.style={draw, align=center, inner sep=5pt, rounded corners=1pt},
    sees/.style={->, >=Stealth, semithick},
    gap/.style={->, >=Stealth, semithick, dashed},
    lbl/.style={font=\footnotesize, align=center},
]
\node[box, minimum width=4.0cm] (hh) at (0,0) {\textbf{Household}\\[1pt]
    \footnotesize income, savings, traditional debt,\\
    \footnotesize \ac{BNPL} obligations, informal debt};
\node[box] (lender) at (0,3.0) {\textbf{Traditional lender}\\ \footnotesize Regulation 23A gate};
\node[box] (bureau) at (7.0,3.0) {\textbf{Credit bureau}\\ \footnotesize traditional debt, defaults};
\node[box] (p1) at (6.0,0.35)  {\footnotesize Platform 1};
\node[box] (p2) at (8.2,0.35)  {\footnotesize Platform 2};
\node[box] (p3) at (6.0,-0.75) {\footnotesize Platform 3};
\node[box] (p4) at (8.2,-0.75) {\footnotesize Platform 4};
\node[draw, dashed, inner sep=6pt, fit=(p1)(p2)(p3)(p4),
      label={[lbl]below:{mutually blind (Submodel 13)}}] (plats) {};
\node[box] (group) at (0,-2.6) {\textbf{Reference group}\\ \footnotesize income quintile $\times$ province};

\draw[sees] (hh) -- (lender) node[lbl, midway, left] {applications:\\ declared income};
\draw[sees, <->] (lender) -- (bureau) node[lbl, midway, above] {reports and reads};
\draw[sees] (hh.east) -- (plats.west) node[lbl, midway, above] {pay-in-four\\ facilities};
\draw[gap] (plats.north) -- (bureau.south) node[lbl, midway, right] {BNPL unreported\\ (Scenario 1 adds visibility)};
\draw[sees] (group) -- (hh) node[lbl, midway, left] {lagged adoption\\ share $s_g$};
\end{tikzpicture}
\caption{Information structure of the model}
\label{fig:visibility}
\tabnote{Boxes represent the five entity types. Solid arrows show the available information flows; dashed connections identify the benchmark's reporting gaps. Platforms cannot observe competitors' exposures (Submodel~13). The bureau excludes \ac{BNPL} unless the reporting switch in Submodel~15 is enabled. The traditional lender's affordability assessment therefore omits these obligations in the benchmark. Households observe their own accounts and the lagged adoption share $s_g$ of their reference group (Submodel~11).}
\end{figure}

The traditional lender assesses applications using declared gross income, its
existing exposures and the traditional balances recorded at the bureau.
In the benchmark, this assessment excludes \ac{BNPL} commitments
\cite{nortonrose_bnpl_sa, transunion_cps_sa_2025} and informal borrowing
\cite{finscope2019}. The model therefore applies an affordability rule to an
incomplete view of household obligations.

The boolean parameter \texttt{bnpl\_bureau\_visible}, written
$\mathbb{1}^{\text{bur}}$, controls whether active \ac{BNPL} obligations enter
this calculation. It is false in the benchmark. Switching it on adds those
obligations to the traditional lender's affordability assessment, representing
one channel of the reporting arrangements discussed in
Section~\ref{sec:background}. A second switch, $\mathbb{1}^{\text{aff}}$,
applies the same affordability gate to \ac{BNPL} originations.
Algorithms~\ref{alg:gate} and~\ref{alg:bnpl} specify how the two switches
enter lending decisions.

The bureau switch can change approval through the affordability calculation,
but it creates no credit score, changes no lending prices and cannot represent
score-based refusals. The omitted scoring channel may be material, given the
preliminary score reductions reported by FINASA \cite{finasa_bnpl_2026}. It is
excluded from the comparison in Section~\ref{sec:scenarios} and discussed as an
extension in Section~\ref{sec:future-work}.

\subsection{Within-Tick Sequence and State Changes}
\label{sec:sequence}

Each simulation step, or tick, represents 14 days. This places the biweekly
pay-in-four instalments directly on tick boundaries. Monthly survey flows are
converted to tick amounts using $12/26$, the ratio of months to ticks in a year. Households act in a newly randomised
order each tick, with credit applications decided as they occur. The peer-adoption
signal is updated once all households have acted.

Each household follows seven steps (Algorithm~\ref{alg:tick}):
\begin{enumerate}
    \item Receive income, adjusted for employment shocks.
    \item Pay committed expenditure from income and savings.
    \item Accrue interest and calculate traditional and \ac{BNPL} payments due, including arrears. Minimum-payers owe the larger of interest and the minimum balance fraction; scheduled-payers owe their contractual instalment.
    \item Seek credit for a shortfall, using \ac{BNPL} first where eligible and traditional credit for the remainder. Assess distress before payment friction.
    \item Collect \ac{BNPL} payments, then pay traditional debt subject to payment friction.
    \item Spend remaining cash on discretionary consumption.
    \item Draw any want-driven \ac{BNPL} purchase and update balances, arrears, distress and default.
\end{enumerate}

The order determines how a shortfall develops. Because committed expenditure is
paid first, essential living costs can leave too little cash for debt service,
consistent with a cash-flow approach to distress \cite{madeira2018chile}.
Credit applications precede repayment, so borrowing can cover that gap
\cite{dorazio2017micro}. Distress is assessed before payment friction,
so skipping an otherwise affordable instalment does not itself classify a
household as unable to pay. Using the previous tick's peer signal also prevents
adoption decisions from depending on which neighbour acts first.

Distress and arrears are distinct states. A household is distressed if it cannot
cover committed expenditure from opening cash and income, or cannot meet
scheduled debt service after seeking credit. An initial committed-expenditure
shortfall remains recorded as distress even if later borrowing supplies enough
cash to cover it. Arrears track unpaid debt obligations through the \ac{CCMR}
age bands. Default occurs after seven consecutive distressed ticks, or 98 days,
and blocks all further credit. The duration is chosen to approximate the
90-day impairment convention at a 14-day resolution; the distress rule is not
identical to an account being 90 days in arrears.

\subsection{Credit and Adoption Mechanisms}
\label{sec:submodels}

The initial population is constructed as described in
Section~\ref{sec:population}.
Table~\ref{tab:submodels} sets out the seventeen operational submodels.
The following subsections explain borrowing amounts, peer adoption and
cross-platform stacking in more detail.

\begingroup
\setlength{\tabcolsep}{4pt}
\small
\begin{longtable}{L{2.35cm}L{4.75cm}L{2.3cm}L{4.0cm}}
\caption{Submodel specification}
\label{tab:submodels} \\
\multicolumn{4}{p{13.4cm}}{\footnotesize The seventeen submodels, each with its rule, source and parameters. Behavioural rules are either grounded in literature or flagged as assumptions, and free parameters are sourced, derived or swept, as defined in the caption of Table~\ref{tab:param-register}, which gives each parameter's default value, provenance and sweep range. Bold marks the rules with no literature anchor, for which sensitivity analysis is conducted. Submodels 4, 5, 10, 11 and 13 are elaborated in the text and Submodel 15 in Section~\ref{sec:scenarios}; the rest are given in full as pseudocode in Appendix~\ref{app:pseudo-main}.} \\
\toprule
\textbf{Submodel} & \textbf{Rule} & \textbf{Source} & \textbf{Parameters and validation} \\
\midrule
\endfirsthead
\toprule
\textbf{Submodel} & \textbf{Rule} & \textbf{Source} & \textbf{Parameters and validation} \\
\midrule
\endhead
\bottomrule
\endfoot

1. Income and shock & Baseline income is static. Each employed member of a wage-dominant household separates with a per-tick probability, and the household loses that member's share of its wage component until re-employment. Grant and remittance-dominant households are exempt: grants are statutory transfers. & \cite{madeira2018chile, statssa_lmd_2022} & Income-shock probability fitted to the \ac{CCMR} 90+ band and reported against the \ac{QLFS} 2017 separation band of 0.54--1.17\% per tick. No magnitude parameter: the loss is the observed wage component. Exit hazard 1.87\% per tick from \ac{QLFS} 2017 Q3--Q4 flows. Single-tick non-persistent shock kept as a robustness arm. \\
\addlinespace
2. Consumption & Two tiers. Committed expenditure (food, rent) is an incompressible floor; discretionary spending compresses to zero before any instalment is missed. & Data layer (Section~\ref{sec:population}) & Default floor of zero; habit-persistence floors of 0.25 and 0.50 swept. \\
\addlinespace
3. Borrowing trigger & Traditional credit is sought only for realised shortfalls. \ac{BNPL} adds a want-driven path governed by peer adoption. & \cite{dorazio2017micro, ackert2025bnpl, dehaan2024bnpl} & $q_{\text{base}}$, $\beta$; see Submodel 11. \\
\addlinespace
4. Borrowing amount & Shortfall path: \textbf{structural assumption} with no anchor located; households request the exact shortfall. Want-driven path: purchase size is drawn lognormally with mean $\kappa$ times the household's own monthly discretionary budget, so purchase size inherits the population's heterogeneity. & \cite{statssa_ies_2023} & Sweeps: shortfall / $+25\%$ / $+$ one tick of committed expenditure; $\kappa = 0.139$ (derived from IES 2022/23 COICOP shares) swept 0.07--0.28; dispersion 0.6 (assumption); sizing base swept discretionary / income. Realised mean order checked against the R992 (2017 Rands) issuer-disclosed average order value \cite{weaver2025iar}. \\
\addlinespace
5. Lender choice & Eligible households take \ac{BNPL} first, up to $\kappa$ times their monthly discretionary budget, an assumed ceiling on its contribution to the shortfall. The remainder routes to traditional credit. The no-\ac{BNPL} baseline uses traditional credit only. & \cite{ackert2025bnpl} & The cap is an \textbf{assumption}; an uncapped arm is crossed with the three Submodel~4 rules. Comparison with Pattern 1. \\
\addlinespace
6. Repayment and arrears & Scheduled-payers pay the full amortised instalment. Minimum-payers pay the larger of accrued interest and 5\% of balance per month, tick-scaled; the fraction is an \textbf{assumption}, since neither the \ac{NCA} nor any located source prescribes a formula. Unpaid amounts accrue as arrears in \ac{CCMR} age bands. Payment friction skips an affordable traditional instalment with fixed probability. & \cite{Keys2019, Kuchler2021} & Minimum-payer share 0.29 (swept 0.20--0.40); minimum fraction swept 0.025--0.10. Friction is the second fitted parameter, calibrated to the \ac{CCMR} 1--30 day band; it applies to traditional debt only, since \ac{BNPL} auto-debits a card. \\
\addlinespace
7. Distress and default & Distress records an unpaid committed-expenditure shortfall before borrowing, or unmet scheduled service after borrowing, assessed before payment friction. Seven consecutive distressed ticks trigger default; a defaulted household is refused all further credit. & \cite{madeira2018chile} & $k=7$ ticks (98 days), the first tick boundary at or past the 90-day impairment convention; $k=4$ (60+ band) swept. \\
\addlinespace
8. Heterogeneity & Behavioural parameters vary by income quintile, demographics, and the financial-inclusion markers of the data layer. & \cite{Keys2019} & No continuous present-bias parameter: no local calibration data. \\
\addlinespace
9. Credit-granting gate & The Regulation 23A residual-income test, applied at household level; all-or-nothing. New credit is priced as unsecured (repo $+21\%$, 25-month term), matching Section~\ref{sec:servicing}. A granted loan is booked onto the household's consolidated balance, raising the balance, its balance-weighted rate and scheduled service, so each grant reduces the headroom left for the next. & \cite{sa_ncr_affordability_2015} & No free parameter; the statutory expense table is hard-coded and asserted against the regulation's published worked examples. \\
\addlinespace
10. Information asymmetry & The bureau records traditional debt and defaults only. Lenders cannot observe \ac{BNPL} or informal debt. & \cite{nortonrose_bnpl_sa, transunion_cps_sa_2025} & \texttt{bnpl\_bureau\_visible} boolean (default off, the pre-2026 position); Scenario 1 in Submodel 15. \\
\addlinespace
11. Peer influence & $q_i(t) = \mathrm{clip}(q_{\text{base}} + \beta s_{g(i)}(t-1), 0, 1)$ within reference groups, where $s_g$ is measured on the group's \ac{BNPL}-eligible members. A heterogeneous-threshold variant is pre-registered as the structural alternative and run on identical grids. & \cite{ackert2025bnpl, cardaci2018inequality, granovetter1978threshold} & $\beta=0$ control arm reported in all comparisons; the threshold arm reports its own $\gamma=0$ row. \\
\addlinespace
12. \ac{BNPL} eligibility & Banked households only; a light automated screen in place of the Regulation 23A test. Per-order cap R9{,}927 (2017 Rands); rolling limit $\lambda = 0.10$ of monthly income per platform, set for each household in line with the industry's disclosed ``low and grow'' practice. & \cite{payflex_terms, payjustnow_terms, homechoice2024ir} & Access capped at the imputed 82.8\% banked share. The default of 0.10 is set so that four per-platform limits sum to the stacked total Woolard observed, roughly 37\% of monthly income across providers; $\lambda$ is swept 0.1--1.0 and reported against a 0.10--0.26 band converted by the author from UK and Australian disclosures \cite{woolard2021, asic2020rep672}. Order-cap and limit binding rates reported by quintile. \\
\addlinespace
13. Stacking & Concurrent facilities across platforms, with no aggregate limit. Platforms stay blind to competitor exposures. & Follows Submodel 10 & Platforms swept 1--6 (default 4); emergent depth compared with the \ac{CFPB} share in Section~\ref{sec:data-validation}. \\
\addlinespace
14. \ac{BNPL} repayment & Exogenous ``Pay-in-4'': 25\% at checkout, then 25\% at each of the next three ticks, interest-free. The checkout debit must clear, so an order is truncated to what the household can fund and liquidity caps purchase size endogenously. Late fees of R125.74 per tick, capped at R188.61 per facility (2017 Rands); a household that exhausts the cap is cut off by that platform only. & \cite{payflex_terms} & No free parameters; instalments fall on tick boundaries. Pay-in-3 monthly (PayJustNow) is the structural alternative arm. \\
\addlinespace
15. Scenario switches & Bureau visibility ($\mathbb{1}^{\text{bur}}$: \ac{BNPL} obligations enter the Regulation 23A test), mandatory affordability screening ($\mathbb{1}^{\text{aff}}$: the Regulation 23A test applied to \ac{BNPL} originations), and peer influence ($\beta > 0$), crossed (Section~\ref{sec:scenarios}). & \cite{businessday_bnpl_2026, fca_ps26_1, woolard2021} & The two regulatory switches are binary; each is compared with the benchmark at $\beta = 0$ and $\beta = 1$. \\
\addlinespace
16. Macro environment & Static. The repurchase rate is fixed at 7.00\%, freezing statutory maximum interest rates. & Experimental design (Section~\ref{sec:purpose}) & Exogenous constraint. \\
\addlinespace
17. Activation order & Random, re-drawn each tick; applications decided inline as each household acts. & \cite{comer2013activation, alizadeh2015activation} & Fixed-order robustness arm (Appendix~\ref{app:supplementary}). \\
\end{longtable}
\endgroup

\subsubsection*{Submodels 4 and 5: Borrowing on a Shortfall}
\label{sec:amount}

When a household faces a shortfall, Submodel~4 determines how much it requests.
By default, the request equals the shortfall. Two alternatives add either 25\% or one
tick of committed expenditure. No located study identifies the appropriate
rule, so all three are compared in the sensitivity analysis
(Section~\ref{sec:results-robustness}).

Submodel~5 determines which lender receives the request. An eligible household
uses \ac{BNPL} first, up to $\kappa$ times its monthly discretionary budget.
The remaining amount is requested from the traditional lender, which grants or
refuses it in full under Submodel~9.

The cap represents a restriction on what retail credit can finance. \ac{BNPL}
is not general-purpose cash, so the model limits its contribution to a shortfall
using the same budget share that sizes want-driven purchases. No empirical
maximum was located, so the sensitivity analysis repeats each borrowing-amount rule with the cap removed.

\subsubsection*{Submodel 11: Peer Influence}
\label{sec:peer-influence}

Peer influence changes the probability of a want-driven purchase. Let
$s_g(t-1)$ be the share of eligible households in reference group $g$ holding
an active \ac{BNPL} balance at the previous tick. Using eligible households
as the denominator separates this signal from the group's overall banking
rate. Household $i$'s purchase propensity is:

\begin{equation}
    q_i(t) = \mathrm{clip}\!\left(q_{\text{base}} + \beta \, s_{g(i)}(t-1),\; 0,\; 1\right).
    \label{eq:peer}
\end{equation}

Here, $\mathrm{clip}$ restricts the probability to $[0,1]$. A positive
$q_{\text{base}}$ allows borrowing to begin before anyone in the group has
adopted. The coefficient $\beta$ determines how strongly the household responds
to others' adoption. Evidence on social approval and consumption emulation
motivates this channel \cite{ackert2025bnpl, cardaci2018inequality} without
identifying its strength, so both parameters are swept, with $\beta=0$
providing the no-peer control.

A Granovetter-style threshold rule provides a structural alternative
\cite{granovetter1978threshold}. Under this rule, the propensity rises to
$q_i = \mathrm{clip}(q_{\text{base}} + \gamma,0,1)$ when the lagged group
share reaches the household's threshold. Thresholds are drawn from
$\mathcal{N}(\mu_\theta, \sigma_\theta)$ and clipped to $[0,1]$, so households
with a negative draw respond immediately.
Fixed initialisation seeds make the control settings comparable. The sweeps
vary threshold dispersion ($\sigma_\theta$), with additional comparisons of
the mean, to examine whether outcomes depend on the distribution of peer
susceptibility. Setting $\gamma=0$ removes the threshold response.

\subsubsection*{Submodel 13: Stacking}
\label{sec:stacking}

Because each platform limits only its own exposure to a household, repeated
borrowing can create concurrent obligations across platforms, with no aggregate
exposure cap in the benchmark. The model does not
assign a target number of providers to each household: cross-platform stacking
arises from borrowing requests and the available platform headroom.

We vary the number of platforms from one to six. One platform rules out
cross-platform stacking, but still permits multiple agreements within that
platform. This distinguishes provider count from the number of concurrent
agreements. The stacking measures are compared with the \ac{CFPB} estimate in
Section~\ref{sec:data-validation}, which also discusses differences in
definitions, markets and observation periods.

\subsection{Implementation and Verification}
\label{sec:implementation}

The model is implemented in Python using Mesa. Households, the lender, the
bureau and platforms have separate classes; reference-group membership is stored
on the model. Household data load from the parquet files described in
Section~\ref{sec:population}. The lending gate reuses the affordability function
used when constructing scheduled debt service, keeping the two calculations
consistent. The parameter register in Appendix~\ref{app:specification} is
generated from the code's declarations.

Automated checks cover balance-sheet accounting, arrears transitions,
Regulation 23A worked examples, pay-in-four repayments and late-fee caps.
Boundary checks disable \ac{BNPL}, remove peer influence ($\beta=0$ or
$\gamma=0$), and restrict access to one platform. All checks pass on the
submitted code.
